%% ============================================================
%%  DKU Signature Work Research Poster — Class of 2026
%%  Flood Inundation Mapping From Sentinel-1 SAR Using
%%  HAND-Guided Gating and Uncertainty Quantification
%%
%%  Author : Bouchra Daddaoui | Computation and Design
%%  Mentor : Prof. Dongmian Zou, Ph.D.
%%
%%  Compile: xelatex poster.tex   (or lualatex)
%%           Works on Overleaf — set compiler to XeLaTeX
%% ============================================================

\documentclass[final]{beamer}

%% ---- Page geometry (A0 portrait) -------------------------
\usepackage[
  orientation=portrait,
  size=a0,
  scale=1.15
]{beamerposter}

%% ---- Core packages ----------------------------------------
\usepackage{lmodern}
\usepackage[T1]{fontenc}
\usepackage[utf8]{inputenc}
\usepackage{microtype}
\usepackage{amsmath,amssymb}
\usepackage{booktabs}
\usepackage{array}
\usepackage{multirow}
\usepackage{xcolor}
\usepackage{graphicx}
\usepackage{tikz}
\usepackage{tcolorbox}
\tcbuselibrary{skins,breakable,most}
\usepackage{enumitem}
\usepackage{multicol}
\usepackage{ragged2e}
\usepackage{setspace}
\usepackage{calc}
\usepackage{url}

%% ---- DKU Colours -----------------------------------------
\definecolor{dkuTeal}{HTML}{4BAAA5}      % primary teal
\definecolor{dkuDarkTeal}{HTML}{2E7873}  % darker teal for rules
\definecolor{dkuLightTeal}{HTML}{D6EEEC} % very light teal background
\definecolor{dkuGold}{HTML}{C9A227}      % accent gold
\definecolor{dkuGrey}{HTML}{4A4A4A}      % body text grey
\definecolor{dkuLightGrey}{HTML}{F5F5F5} % card background
\definecolor{alertRed}{HTML}{C0392B}
\definecolor{successGreen}{HTML}{27AE60}
\definecolor{highlightBlue}{HTML}{2980B9}

%% ---- Remove Beamer navigation ---------------------------
\setbeamertemplate{navigation symbols}{}
\setbeamertemplate{headline}{}
\setbeamertemplate{footline}{}
\setbeamercolor{background canvas}{bg=white}

%% ---- tcolorbox styles ------------------------------------
\tcbset{
  posterbox/.style={
    enhanced,
    boxrule=0pt,
    frame hidden,
    colback=dkuLightGrey,
    borderline west={4pt}{0pt}{dkuTeal},
    arc=4pt,
    left=12pt, right=12pt, top=10pt, bottom=10pt,
    before upper={\setlength{\parskip}{4pt}\justifying\small},
  },
  sectiontitle/.style={
    enhanced,
    boxrule=0pt,
    frame hidden,
    colback=dkuTeal,
    coltext=white,
    arc=4pt,
    left=14pt, right=14pt, top=8pt, bottom=8pt,
    fontupper=\Large\bfseries,
  },
  resultbox/.style={
    enhanced,
    boxrule=1.5pt,
    colframe=dkuTeal,
    colback=white,
    arc=6pt,
    left=10pt, right=10pt, top=8pt, bottom=8pt,
    fontupper=\footnotesize,
  },
  keymetric/.style={
    enhanced,
    boxrule=0pt,
    frame hidden,
    colback=dkuTeal,
    coltext=white,
    arc=8pt,
    halign=center,
    valign=center,
    left=6pt, right=6pt, top=6pt, bottom=6pt,
  },
  keymetric_alt/.style={
    enhanced,
    boxrule=2pt,
    colframe=dkuTeal,
    colback=white,
    coltext=dkuDarkTeal,
    arc=8pt,
    halign=center,
    valign=center,
    left=6pt, right=6pt, top=6pt, bottom=6pt,
  },
  warnbox/.style={
    enhanced,
    boxrule=0pt,
    frame hidden,
    colback=dkuGold!20,
    borderline west={4pt}{0pt}{dkuGold},
    arc=4pt,
    left=12pt, right=12pt, top=8pt, bottom=8pt,
    before upper={\setlength{\parskip}{4pt}\justifying\small},
  },
}

%% ===========================================================
\begin{document}
\begin{frame}[t,fragile]{}

%% ===========================================================
%%  HEADER BLOCK
%% ===========================================================
\begin{tikzpicture}[remember picture, overlay]
  %% Full-width teal banner
  \fill[dkuTeal] (current page.north west)
    rectangle ([yshift=-12.5cm]current page.north east);
  %% Bottom rule on banner
  \fill[dkuDarkTeal] ([yshift=-12.5cm]current page.north west)
    rectangle ([yshift=-13.1cm]current page.north east);
\end{tikzpicture}

\vspace{0.2cm}

%% --- Header content ----------------------------------------
\begin{columns}[T]
  \begin{column}{0.13\textwidth}
    \vspace{0.6cm}
    \hspace{0.8cm}
    %% DKU Logo placeholder — replace with actual DKU logo file
    \begin{tikzpicture}
      \draw[white, line width=3pt, rounded corners=6pt]
        (0,0) rectangle (4.8,4.8);
      \node[white, font=\Large\bfseries] at (2.4,2.8) {DKU};
      \node[white, font=\small] at (2.4,1.8) {Duke Kunshan};
      \node[white, font=\small] at (2.4,1.1) {University};
    \end{tikzpicture}
    %% To use the real logo, comment out the tikzpicture above and uncomment:
    %% \includegraphics[width=4.5cm]{dku_logo_white.png}
  \end{column}

  \begin{column}{0.74\textwidth}
    \vspace{0.5cm}
    \begin{center}
      {\color{white}\fontsize{42}{50}\bfseries\selectfont
       Flood Inundation Mapping From Sentinel-1 SAR\\[6pt]
       Using HAND-Guided Gating and\\[6pt]
       Uncertainty Quantification}\\[14pt]
      {\color{white!90}\fontsize{24}{30}\selectfont
       \textbf{Bouchra Daddaoui} \;|\; Computation and Design \quad
       $\cdot$ \quad
       Mentor: \textbf{Prof. Dongmian Zou, Ph.D.}}\\[6pt]
      {\color{white!80}\fontsize{20}{24}\selectfont
       Duke Kunshan University \;|\; Signature Work \;|\; May 2026}
    \end{center}
  \end{column}

  \begin{column}{0.13\textwidth}
    \vspace{0.6cm}
    \hspace{0.3cm}
    \begin{tikzpicture}
      \fill[white, rounded corners=6pt] (0,0) rectangle (4.8,4.8);
      \node[dkuTeal, font=\fontsize{11}{13}\bfseries, align=center]
        at (2.4,3.2) {Signature\\Work};
      \node[dkuTeal, font=\fontsize{14}{16}\bfseries, align=center]
        at (2.4,2.0) {Class of\\2026};
      \node[dkuTeal, font=\normalsize, align=center]
        at (2.4,0.9) {\textbf{DKU}};
    \end{tikzpicture}
  \end{column}
\end{columns}

\vspace{1.2cm}

%% ===========================================================
%%  KEY METRICS BAR
%% ===========================================================
\hspace*{1.8cm}
\begin{minipage}{0.94\textwidth}
\begin{tcolorbox}[
  enhanced, colback=dkuLightTeal, colframe=dkuTeal,
  boxrule=1.5pt, arc=8pt,
  left=14pt, right=14pt, top=10pt, bottom=10pt,
]
\begin{columns}[c]
  \begin{column}{0.19\textwidth}
    \centering
    {\color{dkuDarkTeal}\fontsize{40}{44}\bfseries\selectfont 0.724}\\
    {\color{dkuGrey}\normalsize Best Model IoU (D$_\text{full}$)}
  \end{column}
  \begin{column}{0.005\textwidth}
    \color{dkuTeal}\rule{1.5pt}{5cm}
  \end{column}
  \begin{column}{0.19\textwidth}
    \centering
    {\color{dkuDarkTeal}\fontsize{40}{44}\bfseries\selectfont 78.6\%}\\
    {\color{dkuGrey}\normalsize ECE Reduction via Cal.}
  \end{column}
  \begin{column}{0.005\textwidth}
    \color{dkuTeal}\rule{1.5pt}{5cm}
  \end{column}
  \begin{column}{0.19\textwidth}
    \centering
    {\color{dkuDarkTeal}\fontsize{40}{44}\bfseries\selectfont 7.84M}\\
    {\color{dkuGrey}\normalsize People at Flood Risk}
  \end{column}
  \begin{column}{0.005\textwidth}
    \color{dkuTeal}\rule{1.5pt}{5cm}
  \end{column}
  \begin{column}{0.19\textwidth}
    \centering
    {\color{dkuDarkTeal}\fontsize{40}{44}\bfseries\selectfont +61.4\%}\\
    {\color{dkuGrey}\normalsize TTA Uncertainty Corr.}
  \end{column}
  \begin{column}{0.005\textwidth}
    \color{dkuTeal}\rule{1.5pt}{5cm}
  \end{column}
  \begin{column}{0.19\textwidth}
    \centering
    {\color{dkuDarkTeal}\fontsize{40}{44}\bfseries\selectfont 446}\\
    {\color{dkuGrey}\normalsize Sen1Floods11 Chips}
  \end{column}
\end{columns}
\end{tcolorbox}
\end{minipage}

\vspace{1.2cm}

%% ===========================================================
%%  MAIN CONTENT — THREE COLUMNS
%% ===========================================================
\begin{columns}[T]

%%------------------------------------------------------------
%%  COLUMN 1 — Abstract + Data
%%------------------------------------------------------------
\begin{column}{0.315\textwidth}
\hspace*{1.8cm}
\begin{minipage}{0.93\textwidth}

%% -- Abstract -----------------------------------------------
\begin{tcolorbox}[sectiontitle]
  \vphantom{Ag} Abstract
\end{tcolorbox}
\vspace{4pt}
\begin{tcolorbox}[posterbox]
Accurate, near-real-time flood mapping is critical for humanitarian response yet
remains challenging due to persistent cloud cover and the spectral ambiguity of SAR
imagery. This Signature Work develops \textbf{TerrainFlood-UQ}: a physics-informed
Siamese ResNet-34 encoder–decoder that (i) integrates Height Above Nearest Drainage
(HAND) topography through a learnable attention gate, and (ii) quantifies prediction
uncertainty via Monte Carlo Dropout and Test-Time Augmentation.

Training on the Sen1Floods11 benchmark (Sentinel-1 SAR; 11 countries) and
\textbf{evaluated out-of-distribution on Bolivia}, the best model achieves an
\textbf{IoU of 0.724}—compared to 0.582 for classical Otsu thresholding and 0.421
for a standard U-Net baseline.  Calibration is dramatically improved by temperature
scaling (ECE: $0.363 \to 0.063$). TTA-based uncertainty reliably flags uncertain
flood boundaries ($r = +0.614$), enabling probabilistic exposure estimates of
\textbf{7.84 million people at flood risk} in Bolivia with 1.21 million (15.4\%)
in high-uncertainty zones.
\end{tcolorbox}

\vspace{14pt}

%% -- Research Gap -------------------------------------------
\begin{tcolorbox}[sectiontitle]
  \vphantom{Ag} Motivation
\end{tcolorbox}
\vspace{4pt}
\begin{tcolorbox}[posterbox]
\begin{itemize}[leftmargin=1.2em,itemsep=3pt]
  \item[\color{dkuTeal}\textbullet]
    Floods affect $\sim$250 million people/year; SAR images through cloud cover but
    water–sand/ice confusion causes high false positives in flat terrain.
  \item[\color{dkuTeal}\textbullet]
    Classical methods (Otsu thresholding) are fast but fragile; deep models
    generalise better but produce \textit{overconfident} predictions without
    uncertainty estimates.
  \item[\color{dkuTeal}\textbullet]
    HAND encodes the \textit{physical likelihood} of inundation: low-lying valleys
    ($h < 30$ m) flood first.  Injecting this prior into the network gate should
    reduce false alarms on elevated areas.
  \item[\color{dkuTeal}\textbullet]
    No prior work combines HAND gating \textit{and} calibrated UQ for operational
    SAR flood mapping.
\end{itemize}
\end{tcolorbox}

\vspace{14pt}

%% -- Dataset ------------------------------------------------
\begin{tcolorbox}[sectiontitle]
  \vphantom{Ag} Dataset — Sen1Floods11
\end{tcolorbox}
\vspace{4pt}
\begin{tcolorbox}[posterbox]
\textbf{Satellite data:} Sentinel-1 C-band SAR (VV + VH polarisations)
collected pre- and post-flood; 446 hand-labelled $512 \times 512$ chips across
11 countries covering diverse flood typologies.\\[4pt]

\textbf{Auxiliary:} HAND raster from MERIT Hydro DEM (z-score normalised in the
data tensor; denormalised to metres inside the model).\\[4pt]

\textbf{Splits:}
\begin{itemize}[leftmargin=1.2em,itemsep=2pt,topsep=2pt]
  \item Train / Val: 10 countries; stratified by flood fraction
  \item \textbf{Test (OOD): Bolivia only} — 15 chips, 2,867,815 valid pixels
    \\ \textit{(Bolivia withheld throughout; never seen during training)}
\end{itemize}

\vspace{4pt}
\begin{tabular}{@{}lrr@{}}
  \toprule
  Split & Chips & Flood px \\
  \midrule
  Train & 340 & 22.8\% \\
  Val   & 91  & 21.6\% \\
  Test (Bolivia) & 15 & 18.3\% \\
  \bottomrule
\end{tabular}

\vspace{6pt}
Each $512\times512$ chip at 10\,m GSD covers ${\approx}26\,\text{km}^2$.
Six bands are stacked per chip:
VV$_\text{pre}$, VH$_\text{pre}$, VV$_\text{post}$, VH$_\text{post}$,
VV/VH ratio (unused), HAND.
\end{tcolorbox}

\vspace{14pt}

%% -- Why Bolivia --------------------------------------------
\begin{tcolorbox}[warnbox]
\textbf{Why Bolivia as holdout?} Bolivia's 2018 Amazonian flood was catastrophic
(80,000 displaced) but is climatically \textit{unlike} the training countries,
testing true out-of-distribution generalisation—the hardest and most realistic
evaluation scenario.
\end{tcolorbox}

\end{minipage}
\end{column}

%%------------------------------------------------------------
%%  COLUMN 2 — Methods
%%------------------------------------------------------------
\begin{column}{0.37\textwidth}

%% -- Architecture -------------------------------------------
\begin{tcolorbox}[sectiontitle]
  \vphantom{Ag} Model Architecture — TerrainFlood-UQ
\end{tcolorbox}
\vspace{4pt}
\begin{tcolorbox}[posterbox]

%% Architecture diagram (ASCII-art schematic)
\begin{center}
\begin{tikzpicture}[
  scale=0.88, every node/.style={font=\small},
  box/.style={draw=dkuTeal, fill=dkuLightTeal, rounded corners=4pt,
              minimum width=2.6cm, minimum height=0.75cm, align=center},
  smallbox/.style={draw=dkuGold, fill=dkuGold!20, rounded corners=3pt,
                   minimum width=2.2cm, minimum height=0.6cm, align=center,
                   font=\footnotesize},
  arrow/.style={->, thick, dkuDarkTeal},
]
  %% Pre branch
  \node[box] (pre) at (0, 4) {Pre-event\\VV, VH};
  \node[box] (enc1) at (0, 2.4) {ResNet-34\\Encoder};
  \draw[arrow] (pre) -- (enc1);

  %% Post branch
  \node[box] (post) at (4.5, 4) {Post-event\\VV, VH};
  \node[box] (enc2) at (4.5, 2.4) {ResNet-34\\Encoder};
  \draw[arrow] (post) -- (enc2);

  %% Weight-share annotation
  \draw[dashed, dkuGrey, thick] (enc1.north east) -- (enc2.north west)
    node[midway, above, font=\scriptsize, dkuGrey]{shared weights};

  %% HAND branch
  \node[box, fill=dkuGold!20, draw=dkuGold] (hand) at (9, 4)
    {HAND\\(metres)};
  \node[smallbox] (gate) at (9, 2.4) {Attention\\Gate $\alpha$};
  \draw[arrow] (hand) -- (gate);
  \node[font=\scriptsize, dkuGrey, align=center] at (9, 1.7)
    {$\alpha = \exp(-h/50)$};

  %% Fusion
  \node[box, minimum width=3.0cm] (fuse) at (4.5, 0.9) {Feature\\Fusion};
  \draw[arrow] (enc1.south) |- (fuse.west);
  \draw[arrow] (enc2.south) -- (fuse.north);
  \draw[arrow] (gate.south) |- (fuse.east);

  %% Decoder
  \node[box, minimum width=3.0cm] (dec) at (4.5, -0.6) {U-Net\\Decoder};
  \draw[arrow] (fuse) -- (dec);

  %% MC Dropout annotation
  \node[smallbox] (mcd) at (9, -0.6) {MC Dropout\\(T=20)};
  \draw[dashed, arrow, dkuGold] (mcd) -- (dec);

  %% Output
  \node[box, fill=dkuTeal, draw=dkuDarkTeal,
        text=white, minimum width=3.0cm] (out) at (4.5, -2.1)
    {Flood Map\\$+$ Uncertainty};
  \draw[arrow] (dec) -- (out);
\end{tikzpicture}
\end{center}

\vspace{6pt}
\textbf{Key design choices:}
\begin{itemize}[leftmargin=1.2em,itemsep=3pt]
  \item \textbf{Siamese encoder}: Weight-shared ResNet-34 processes pre/post SAR
    images independently, producing change-aware feature maps.
  \item \textbf{HAND attention gate}: Physical gating $\alpha = \exp(-h/50)$
    suppresses elevated terrain features; $\alpha \approx 0$ at $h > 150$ m.
  \item \textbf{MC Dropout (Variant D)}: Dropout layers remain \textit{active at
    inference}; T=20 stochastic forward passes estimate epistemic uncertainty as
    predictive variance.
  \item \textbf{Temperature scaling}: Post-hoc calibration with learned scalar $T$
    applied to logits before softmax.
\end{itemize}
\end{tcolorbox}

\vspace{10pt}

%% -- Ablation Variants --------------------------------------
\begin{tcolorbox}[sectiontitle]
  \vphantom{Ag} Ablation Study — Four Variants
\end{tcolorbox}
\vspace{4pt}
\begin{tcolorbox}[posterbox]
\begin{center}
\begin{tabular}{@{}clll@{}}
  \toprule
  \textbf{Var.} & \textbf{Encoder input} & \textbf{HAND use} & \textbf{UQ} \\
  \midrule
  A & VV, VH (2ch) & None & — \\
  B & VV, VH, HAND (3ch) & Concatenated & — \\
  C & VV, VH (2ch) & Attention gate & — \\
  \textbf{D} & \textbf{VV, VH (2ch)} & \textbf{Attention gate} & \textbf{MC Dropout} \\
  \bottomrule
\end{tabular}
\end{center}

\vspace{6pt}
Variants C$_\text{full}$ and D$_\text{full}$ are fully-retrained with 120 epochs
and early stopping (patience 25), replacing the 50-epoch checkpoints.\\[4pt]

Variant D+ adds dropout to encoder \textit{and} decoder (degenerate);
Variant E uses change-difference encoding only (degenerate).
\end{tcolorbox}

\vspace{10pt}

%% -- Uncertainty Quantification -----------------------------
\begin{tcolorbox}[sectiontitle]
  \vphantom{Ag} Uncertainty Quantification
\end{tcolorbox}
\vspace{4pt}
\begin{tcolorbox}[posterbox]

\textbf{MC Dropout} (epistemic):
$$\hat{\sigma}^2_\text{MC} = \frac{1}{T}\sum_{t=1}^{T} p_t^2 - \bar{p}^2, \quad T=20$$

\textbf{Test-Time Augmentation} (aleatoric + epistemic):
$$\hat{\sigma}^2_\text{TTA} = \text{Var}\!\left(\{f_\theta(A_k(\mathbf{x}))\}_{k=1}^{8}\right)$$
using the D$_4$ symmetry group (4 rotations $\times$ 2 flips = 8 augmentations).

\vspace{8pt}
\begin{center}
\begin{tabular}{@{}lcc@{}}
  \toprule
  Method & Mean Variance & Error Corr. $r$ \\
  \midrule
  MC Dropout & $4 \times 10^{-4}$ & $-0.815$ {\color{alertRed}$\downarrow$ inverted} \\
  \textbf{TTA} & $8.3 \times 10^{-3}$ & $\mathbf{+0.614}$ {\color{successGreen}$\uparrow$ useful} \\
  \bottomrule
\end{tabular}
\end{center}

\vspace{6pt}
TTA produces $20-50\times$ higher variance than MC Dropout and a positive
correlation with prediction error—meaning \textit{high uncertainty reliably
flags wrong predictions}. MC Dropout's inverted correlation makes it unsuitable
as a standalone uncertainty estimator for this task.
\end{tcolorbox}

\end{column}

%%------------------------------------------------------------
%%  COLUMN 3 — Results & Conclusions
%%------------------------------------------------------------
\begin{column}{0.315\textwidth}
\begin{minipage}{0.93\textwidth}

%% -- Main Results -------------------------------------------
\begin{tcolorbox}[sectiontitle]
  \vphantom{Ag} Results — Bolivia Test Set (OOD)
\end{tcolorbox}
\vspace{4pt}
\begin{tcolorbox}[resultbox]
\begin{center}
\begin{tabular}{@{}lcccc@{}}
  \toprule
  \textbf{Model} & \textbf{IoU} & \textbf{F1} & \textbf{ECE} & \textbf{CI 95\%} \\
  \midrule
  \color{alertRed} Otsu thresh. & 0.582 & 0.736 & — & — \\
  \color{alertRed} Baseline U-Net & 0.421 & 0.593 & — & — \\
  \midrule
  Variant A & 0.408 & 0.580 & — & [0.24, 0.57] \\
  Variant B & 0.441 & 0.612 & — & — \\
  Variant C & 0.662 & 0.796 & — & [0.45, 0.78] \\
  Variant D & 0.690 & 0.816 & 0.078 & [0.46, 0.77] \\
  \midrule
  C$_\text{full}$ & 0.706 & 0.827 & — & — \\
  \rowcolor{dkuLightTeal}
  \textbf{D$_\text{full}$} & \textbf{0.724} & \textbf{0.840} & \textbf{0.063} & \textbf{[0.50, 0.81]} \\
  \bottomrule
\end{tabular}
\end{center}
\vspace{4pt}
{\footnotesize Bootstrap CIs from 1,000 resamples. ECE after temperature scaling ($T=0.100$).}
\end{tcolorbox}

\vspace{10pt}

%% -- Calibration --------------------------------------------
\begin{tcolorbox}[sectiontitle]
  \vphantom{Ag} Calibration
\end{tcolorbox}
\vspace{4pt}
\begin{tcolorbox}[posterbox]
Temperature scaling ($T = 0.100$, sharpening) reduces
\textbf{ECE from 0.363 to 0.063}—a \textbf{78.6\%} reduction.
Post-calibration reliability curve closely follows the diagonal, confirming
well-calibrated flood probability outputs.\\[6pt]

\textbf{AURC} (selective prediction): $0.5168$\\
\textit{Selective abstention on uncertain pixels improves accuracy, confirming
that uncertainty is meaningful for triage.}
\end{tcolorbox}

\vspace{10pt}

%% -- Population Exposure ------------------------------------
\begin{tcolorbox}[sectiontitle]
  \vphantom{Ag} Population Exposure (Bolivia)
\end{tcolorbox}
\vspace{4pt}
\begin{tcolorbox}[posterbox]
Flood predictions are intersected with WorldPop population rasters
(not a model input; used post-prediction only).

\vspace{6pt}
\begin{center}
\begin{tabular}{@{}lrr@{}}
  \toprule
  \textbf{Zone} & \textbf{People} & \textbf{Share} \\
  \midrule
  Total predicted flood & 7,840,200 & 100\% \\
  High uncertainty ($\tau=0.01$) & 1,210,400 & 15.4\% \\
  Confident flood & 6,629,800 & 84.6\% \\
  \bottomrule
\end{tabular}
\end{center}

\vspace{6pt}
\textbf{Actionable insight:} The 1.21 million people in uncertain zones can be
flagged for priority ground-truth verification or additional SAR acquisitions.
Confident predictions (6.63M) support immediate evacuation routing.
\end{tcolorbox}

\vspace{10pt}

%% -- HAND Gate Analysis -------------------------------------
\begin{tcolorbox}[sectiontitle]
  \vphantom{Ag} HAND Gate Analysis
\end{tcolorbox}
\vspace{4pt}
\begin{tcolorbox}[posterbox]
\begin{itemize}[leftmargin=1.2em,itemsep=3pt]
  \item HAND gating (Variant C vs A): \textbf{+25.5 pp IoU} improvement.
  \item Gate $\alpha = \exp(-h/50)$: near-unity at valley floors ($h < 30$ m),
    suppresses predictions on ridges ($h > 150$ m, $\alpha < 0.05$).
  \item Per-chip analysis: high-HAND chips (elevated terrain) benefit most from
    gating; low-HAND chips already well-captured by SAR alone.
  \item Variant B (HAND concatenated to encoder) underperforms C (gate),
    confirming that multiplicative gating is more effective than additive fusion.
\end{itemize}
\end{tcolorbox}

\vspace{10pt}

%% -- Conclusions --------------------------------------------
\begin{tcolorbox}[sectiontitle]
  \vphantom{Ag} Conclusions
\end{tcolorbox}
\vspace{4pt}
\begin{tcolorbox}[posterbox]
\begin{itemize}[leftmargin=1.2em,itemsep=3pt]
  \item[\color{successGreen}\textbf{1.}]
    \textbf{HAND gating substantially improves flood mapping}: +25 pp IoU over
    SAR-only; the physical prior is more effective as a multiplicative gate than
    as a concatenated input.
  \item[\color{successGreen}\textbf{2.}]
    \textbf{Calibration matters}: Raw deep model is severely overconfident
    (ECE=0.363); temperature scaling reduces to 0.063—suitable for probabilistic
    decision support.
  \item[\color{successGreen}\textbf{3.}]
    \textbf{TTA is the right UQ method for this task}: Variance $20\text{–}50\times$
    MC Dropout, positive error correlation ($r=+0.614$), directly actionable for
    uncertain pixel flagging.
  \item[\color{dkuGold}\textbf{4.}]
    \textbf{Limitations}: MC Dropout produces inverted uncertainty signal;
    model generalises to Bolivia but performance on other OOD regions unknown.
  \item[\color{highlightBlue}\textbf{5.}]
    \textbf{Future work}: Multi-temporal fusion, optical data fusion, end-to-end
    calibration training, sub-national exposure disaggregation.
\end{itemize}
\end{tcolorbox}

\vspace{10pt}

%% -- Acknowledgements ----------------------------------------
\begin{tcolorbox}[
  enhanced, boxrule=0pt, frame hidden,
  colback=dkuLightTeal, arc=4pt,
  left=10pt, right=10pt, top=6pt, bottom=6pt,
  fontupper=\footnotesize,
]
\textbf{Acknowledgements:} The author thanks Prof. Dongmian Zou for mentorship
throughout the Signature Work. Computations used the DKUCC HPC cluster. Data:
Sen1Floods11 (\textit{Bonafilia et al., 2020}), MERIT Hydro HAND
(\textit{Yamazaki et al., 2019}), WorldPop (\textit{Stevens et al., 2015}).
\end{tcolorbox}

\end{minipage}
\end{column}

\end{columns}

\vspace{1.2cm}

%% ===========================================================
%%  FOOTER BAR
%% ===========================================================
\begin{tikzpicture}[remember picture, overlay]
  \fill[dkuTeal] (current page.south west)
    rectangle ([yshift=4.8cm]current page.south east);
  \fill[dkuDarkTeal] ([yshift=4.8cm]current page.south west)
    rectangle ([yshift=5.3cm]current page.south east);
\end{tikzpicture}

\begin{columns}[c]
  \begin{column}{0.35\textwidth}
    \hspace{1.8cm}
    {\color{white}\small
      \textbf{Bouchra Daddaoui} \;|\; Computation and Design, DKU\\[2pt]
      \url{bd101@duke.edu} \;|\; Class of 2026}
  \end{column}
  \begin{column}{0.30\textwidth}
    \centering
    {\color{white}\large\bfseries
      \textbf{Signature Work — Class of 2026}}\\[2pt]
    {\color{white!80}\small Duke Kunshan University}
  \end{column}
  \begin{column}{0.35\textwidth}
    \raggedleft
    \hspace{0pt}
    {\color{white}\small
      TerrainFlood-UQ | Sentinel-1 SAR Flood Segmentation\\[2pt]
      Mentor: \textbf{Prof. Dongmian Zou, Ph.D.}\hspace{1.8cm}}
  \end{column}
\end{columns}

\end{frame}
\end{document}
